%compile then compile with biber (for the biblio) then recompile
\documentclass[12pt]{article}

\usepackage{a4wide} % increase the typeset area
\usepackage{bm}
\usepackage{amsmath,amssymb}
\usepackage{enumitem}
\usepackage{graphicx}
\usepackage{color}
\usepackage{float} %to place figure H
\usepackage{multirow} %multirow for table
%\usepackage[super]{natbib} %exponant biblio
\usepackage{fourier} %for danger sign
\usepackage{chngcntr}
\usepackage{pifont} %for more symbol in enumerate

% Useful packages for table

\usepackage{array,multirow,makecell}
%\usepackage[linesnumbered]{algorithm2e}
\usepackage[table]{xcolor}
\setcellgapes{1pt}
\makegapedcells
\newcolumntype{R}[1]{>{\raggedleft\arraybackslash }b{#1}}
\newcolumntype{L}[1]{>{\raggedright\arraybackslash }b{#1}}
\newcolumntype{C}[1]{>{\centering\arraybackslash }b{#1}}



%hyperref
\usepackage[colorlinks]{hyperref}
\hypersetup{
	colorlinks=true,
	linkcolor=blue,
	filecolor=magenta,      
	urlcolor=blue,
	citecolor=blue
}

% captions
\usepackage{caption}
\newcommand{\vect}[1]{\hat{\boldsymbol{#1}}}
\usepackage{subcaption}
\counterwithin{figure}{section}
\makeatletter
\usepackage[labelformat=simple]{subcaption}
\newcommand\captionsubfigure{%
	\renewcommand\p@subfigure{}
	\renewcommand\thesubfigure{\thefigure.\alph{subfigure}}
}
\makeatother

%to make the appendix
\usepackage{appendix}

%code
\usepackage{listings}
\definecolor{backcolor}{RGB}{240, 240, 240}
\lstdefinestyle{bash}{
	commentstyle=\color{green},
	morecomment=[l][\color{magenta}]{\#},
	backgroundcolor=\color{backcolor},  
	breakatwhitespace=false,
	keepspaces=true,                        
	showspaces=false,                
	showstringspaces=false,
	showtabs=false,                  
	tabsize=1    
}

%for footnote
\usepackage[symbol]{footmisc}
\renewcommand{\thefootnote}{\fnsymbol{footnote}}

%bibliography (with section)
\usepackage[backend=biber,style=numeric,sorting=nyt]{biblatex}
\addbibresource{biblio_ref.bib}
\addbibresource{biblio_doc.bib}


% Titlepage
\newcommand{\reporttitle}{Internship Report : Innovative non-conformal finite element methods for augmented surgery}
\newcommand{\reportauthorOne}{Frédérique Lecourtier}
\newcommand{\reportsupervisorOne}{Michel Duprez}
\newcommand{\reportsupervisorTwo}{Emmanuel Franck}
\newcommand{\reporttype}{Coursework}

%Algorithm
\usepackage{xcolor}
\usepackage[linesnumbered,noline,boxed,commentsnumbered]{algorithm2e}
%\usepackage[noline, linesnumbered]{algorithm2e}% http://ctan.org/pkg/algorithm2e
\SetNlSty{bfseries}{\color{black}}{}


\begin{document}
	\nocite{*}
	
	\begin{titlepage}

\newcommand{\HRule}{\rule{\linewidth}{0.5mm}}

\begin{center}
	\includegraphics[width = 0.5\linewidth]{images/logo-mimesis.png} \\ [1.5cm] 

	\textsc{\Large University of Strasbourg}\\[0.5cm] 
	\textsc{\large Master CSMI}\\[0.95cm] 
	
	\HRule \\[0.4cm]
	\huge\bfseries\reporttitle\par % Title of your document
	\HRule \\[0.4cm]
\end{center}

\vspace{1cm}

\begin{flushleft} \large
	\begin{minipage}{0.4\hsize}
		\textit{Authors:}\\
		\reportauthorOne
	\end{minipage} \hfill 
	\begin{minipage}{0.4\hsize}
		\textit{Supervisors:}\\
		\reportsupervisorOne\\
		\reportsupervisorTwo
	\end{minipage}
\end{flushleft}
\vspace{3 cm}
\makeatletter
Date: \@date
\hfill
\includegraphics[width = 0.2\linewidth]{images/inria.png}\\[1.5cm] 

\vfill % Fill the rest of the page with whitespace



\makeatother


\end{titlepage}


	
	\tableofcontents
	
	\newpage
	
	\textbf{PLAN :}
	
	\begin{enumerate}
		\item Intro : Ce stage est un stage de M2, dans le cadre du master CSMI. C'est la suite d'un projet efféctué pendant le premier semestre.... explication rapide du projet (colab entre Cemosis et Mimesis dans le but de blablabla)
		\begin{enumerate}
			\item Présentation de Mimesis
			\item Contexte : Mimesis est spécialisé dans la simulation en temps réel de blabla. C'est pourquoi ils ont développé une méthode appelée PhiFEM (explication des intérêts de PhiFEM).
			\item Objectifs : On cherche ici à améliorer la précision de la solution ainsi que les temps pour l'obtenir en passant par le biais d'un FNO que l'on va entraîner avec des solutions PhiFEM qui sont bien adapté à ce type de réseau de neurones, due aux grilles cartésienne.
		\end{enumerate}
		\item On veut expliquer les trucs suivants : \\
		PhiFEM / FEM / FNO / 2 Différentes méthodes de Correction (2 cas tests solution analytique, sol FNO) / Legendre / Modèle Dense ?  \\
		On veut mettre un schéma qui explique Entrainemùent du FNO -> Sortie -> Correction (-> Legendre ?)
		\item Conclu
		\item Bibliography
		\item Appendix (Organisation of the repository, Documentation, Github actions ...)
	\end{enumerate}

	\newpage
	
	
	
	% \input{...}
	
\end{document}